\chapter{Estado de la cuestión}

El presente capítulo ofrece una revisión detallada del panorama actual respecto a las soluciones de navegación espacial existentes en el ámbito de los videojuegos.
A lo largo de las siguientes secciones se explorarán las herramientas con mayor relevancia en el mercado actual.
Posteriormente, a partir de un análisis crítico de estas, se justificará la necesidad de la herramienta de código abierto propuesta en este trabajo. 

\section{Herramientas disponibles}

En la actualidad, existe un amplio repertorio de herramientas orientadas a resolver el problema de la navegación en entornos virtuales.
Entre estas soluciones, aquellas que gozan de mayor relevancia y adopción en la industria, se fundamentan en enfoques que combinan estructuras \acrshort{svo} y el algoritmo de búsqueda heurística \gls{astar}~\cite{Millington}, conceptos ya explicados previamente en la Sección~\ref{sec:representacionespacionavegable}.
A nivel arquitectónico, estas soluciones se dividen en dos grandes grupos: aquellas que operan de forma independiente al motor gráfico y las diseñadas para integrarse en un entorno de desarrollo concreto.

\subsection{Herramientas genéricas}

Si se analizan el conjunto de herramientas que son agnósticas al motor, se observa que la solución predominante en el mercado es la de Mercuna~\cite{Mercuna}.
Esta herramienta fue concebida originalmente como un \textit{plugin} comercial exclusivo para Unreal Engine~\cite{unreal}.
Sin embargo, su éxito motivó su posterior extensión para operar de forma desacoplada a este.

El principal atractivo de Mercuna radica en su capacidad para unificar, bajo un mismo ecosistema, prácticamente la totalidad de los casos de uso que engloba la navegación.
Entre sus características más relevantes se encuentran:

\begin{itemize}
    \item \textbf{Navegación volumétrica:} Permite a los agentes tanto voladores como submarinos, moverse por la totalidad del espacio navegable, haciendo uso completo de sus seis grados de libertad.
    \item \textbf{Navegación sobre superficies:} Proporciona la capacidad de desplazarse sobre cualquier tipo de geometría.
    Esto no se limita al suelo convencional, sino que incluye también paredes o techos, gestionando la gravedad de manera local al agente (ver Figura~\ref{fig:anysurface}).
    \item \textbf{Soporte para agentes \glspl{noholonomico}:} Incorpora las restricciones físicas y cinemáticas de los agentes en la fase de cálculo de trayectorias. De este modo, es posible modelar el comportamiento de vehículos con radios de giro limitados o entidades que carecen de la capacidad de acelerar y frenar de manera instantánea.
\end{itemize}

Se trata, sin lugar a dudas, de la herramienta que se puede considerar más robusta y completa del mercado.
Sin embargo, su modelo de negocio supone una barrera de entrada prohibitiva para la mayor parte de estudios.
Mercuna dispone de distintas licencias que dependerán del subconjunto de características requerido, con costes que parten de \pounds15000 o de \pounds2500 para estudios con un presupuesto de desarrollo inferior a \$500000~\cite{MercunaLicensing}.

\begin{figure}[htp]
  \centering
  \includegraphics[width=0.75\textwidth]{imaxes/surfacenavigation.png}
  \caption[Visualización de navegación por cualquier superficie.]
  {Visualización de navegación por cualquier superficie.
  Imagen extraída de~\cite{Mercuna_AnySurface}.}
  \label{fig:anysurface}
\end{figure}

\subsection{Soluciones específicas de motor}

Por el contrario, las herramientas especializadas en un motor gráfico concreto, operan en un paradigma de desarrollo diferente.
Al prescindir de una arquitectura genérica y agnóstica que asegure su interoperabilidad multiplataforma, estas soluciones pueden aprovechar las rutinas nativas del motor para optimizar el rendimiento computacional, ofreciendo asimismo una integración total con el editor.
Esto se traduce en herramientas significativamente más intuitivas para el desarrollador final.

En este ámbito predomina el ecosistema de Unity, fuertemente respaldado por su tienda oficial de recursos: la Asset Store~\cite{assetstore}.
Esta plataforma permite comercializar tanto herramientas de desarrollo como recursos artísticos (modelos tridimensionales, texturas o composiciones musicales, entre otros) que aceleran los flujos de trabajo al evitar la necesidad de diseñar dichos componentes desde el principio.
Bajo un modelo de reparto de ingresos, Unity proporciona visibilidad de los productos creados.

Como consecuencia de la vasta magnitud de este ecosistema, prácticamente la totalidad de las herramientas de terceros se desarrollan de manera exclusiva para este motor.
Mientras que en entornos como Unreal Engine las soluciones complejas suelen comercializarse con costosas licencias comerciales, la Asset Store de Unity proporciona productos más asequibles para estudios independientes o de menor presupuesto.
Entre estas destacan:

\subsubsection{Nav3D (Softbrix Studio)}

Desarrollada por Softbrix Studio, Nav3D es la opción que ostenta mejores valoraciones en la tienda de Unity (Asset Store)~\cite{Nav3D_SoftbrixAssetStore}.
Su mayor virtud reside en que no se limita al cálculo de rutas estáticas, sino que integra un sistema de evasión local.
Esto garantiza que los \acrshort{npc}, en escenarios multiagente, coordinen sus trayectorias y las reajusten en tiempo real para evitar colisionar entre ellos~\cite{Nav3DPro}.

\subsubsection{Nav3D (m6io)}

Creada por el desarrollador independiente m6io, esta solución se comercializa con un enfoque puramente técnico.
A pesar de cubrir la mayoría de casos de uso, carece por completo de herramientas visuales e interfaces de configuración del editor.
En consecuencia, presenta una curva de aprendizaje notablemente más pronunciada, exigiendo al usuario un nivel avanzado de programación en el lenguaje C\#~\cite{Nav3D_FlyingAIAssetStore}.

\subsubsection{HyperNav}

Elaborada por Infohazard Games, constituye actualmente una de las alternativas más robustas y sofisticadas del mercado de Unity, lo cual se refleja también en un coste superior.
Su principal ventaja competitiva es su enfoque multiparadigma: además de la navegación puramente volumétrica, ofrece navegación sobre superficies, permitiendo a los agentes desplazarse por suelos y paredes con una gestión automatizada de la gravedad local~\cite{HyperNav3D}.

\subsubsection{Spatial Pathfinding}

Desarrollada por TecMaid, se presenta como una alternativa más económica para desarrolladores con presupuestos ajustados.
Se trata de una solución que sacrifica funcionalidades avanzadas en favor de ofrecer un sistema más ligero, directo y de implementación rápida~\cite{SpatialPathfinding3D}.

\section{Comparativa de mercado}

A modo de síntesis, la herramienta de carácter agnóstico Mercuna ofrece una solución tecnológicamente superior, si bien su modelo comercial la orienta exclusivamente a producciones de alto presupuesto conocidas como \gls{aaa}.
Por otro lado, las alternativas nativas para Unity democratizan el acceso a estas tecnologías con precios mucho más asequibles, situándose por lo general por debajo de los 100 euros.

Sin embargo, estas últimas opciones se distribuyen bajo modelos de \textbf{código cerrado}, lo que supone una desventaja técnica importante.
Esta opacidad impide a los desarrolladores consultar el código y prever el impacto real de cada funcionalidad en el rendimiento, lo que habitualmente deriva en un uso subóptimo del sistema.
Además, en muchos casos, estas soluciones carecen de un nivel de refinamiento adecuado en cuanto a usabilidad; como es el caso del ya mencionado anteriormente Nav3D de m6io, que prescinde por completo de herramientas gráficas.
Esta comparativa general queda reflejada de manera resumida en la Tabla~\ref{tab:alternativas}.

\begin{table}[htp]
  \centering
  \rowcolors{2}{white}{udcgray!25}
  \begin{tabular}{c|c|c|c}
  \rowcolor{udcpink!25}
  \textbf{Herramienta} & \textbf{Plataforma} & \textbf{Enfoque} & \textbf{Precio} \\\hline
  Mercuna & Multiplataforma & Multiparadigma & \pounds15000 + \\
  \makecell{Nav3D \\ (Softbrix Studio)} & Unity & Volumétrica & 64,39\euro \\
  Nav3D (m6io) & Unity & Volumétrica & 59,80\euro \\
  HyperNav & Unity & \makecell{Volumétrica + \\ Superficies} & 82,80\euro \\
  Spatial Pathfinding & Unity & Volumétrica & 27,59\euro \\
  \end{tabular}
  \caption{Resumen de las alternativas existentes}
  \label{tab:alternativas}
\end{table}

\subsection{Justificación del entorno de desarrollo}
\label{sec:justificacionentorno}

En base a las soluciones analizadas, la elección de Unity como entorno para de desarrollo resulta óptima para este trabajo.
El ecosistema de la Asset Store garantiza no solo que la solución responda a una demanda real, sino que también asegura un canal con la visibilidad suficiente para su correcta difusión.

A este factor se suma una ventaja técnica fundamental: la disponibilidad de herramientas orientadas al alto rendimiento integradas de forma nativa en el motor, como es el caso del Burst Compiler.
Esta tecnología permite compilar código C\# en código máquina altamente optimizado mediante el uso de la infraestructura \gls{llvm}~\cite{llvm2004}.
Dado que la navegación volumétrica implica una carga computacional intensiva, aprovechar el Burst Compiler facilita la ejecución de estas rutinas matemáticas en \textit{hilos secundarios} con un impacto mínimo en el \textit{hilo principal de procesamiento}, preservando así la tasa de refresco del entorno interactivo.

\section{\texorpdfstring{\acrshort{dafo}}{DAFO}}

Con el propósito de evaluar la viabilidad estratégica y justificar el desarrollo del proyecto, se ha llevado a cabo un análisis \acrshort{dafo}.
Este estudio se estructura en un bloque de \textbf{análisis interno} (orientado a identificar las fortalezas y debilidades inherentes a la propuesta), un bloque de \textbf{análisis externo} (centrado en discernir las oportunidades y amenazas que presenta el mercado actual) y un apartado de conclusiones.

\subsection{Análisis interno}

La principal fortaleza del sistema propuesto radica en su condición de ser la primera herramienta de código abierto específicamente orientada a la navegación volumétrica.
A diferencia de las alternativas comerciales cerradas, este enfoque permite a los usuarios auditar el sistema, adaptarlo mediante modificaciones personalizadas según las necesidades de su proyecto y, potencialmente, contribuir a la evolución de este.
Asimismo, la transparencia en el diseño del algoritmo permite a los desarrolladores prever con precisión el impacto en el rendimiento de cada componente.

Por el contrario, la principal debilidad se asocia a su naturaleza como proyecto emergente dentro de un mercado maduro y consolidado por soluciones veteranas.
Esto puede generar una percepción inicial de falta de madurez o robustez, lo que exigirá un proceso continuo de iteración y refinamiento técnico basado en la retroalimentación recibida tras las fases iniciales de despliegue.

\subsection{Análisis externo}

Las oportunidades del proyecto emergen ante la carencia de soluciones integradas o gratuitas de navegación volumétrica en las versiones nativas de los motores gráficos actuales.
En este contexto, la existencia de un mercado centralizado como la Asset Store, permitirá dar a conocer la herramienta directamente a desarrolladores independientes.

Por último, las amenazas provienen del dinamismo y la constante evolución tecnológica que caracterizan al sector de los videojuegos.
El auge de motores alternativos de código abierto, como Godot~\cite{godot}, los cuales incorporan continuamente nuevas capacidades~\cite{salmela2022game}, representa un factor de competitividad externa.
Asimismo, la eventual inclusión de una herramienta de navegación tridimensional oficial y nativa por parte de la propia compañía Unity constituye el riesgo más significativo para soluciones desarrolladas por terceros.

\begin{table}[htp]
  \centering
  \renewcommand{\arraystretch}{1.5}
  \begin{tabular}{|p{0.46\textwidth}|p{0.46\textwidth}|}
  \hline
  \rowcolor{udcpink!25}
  \textbf{Fortalezas} & \textbf{Debilidades} \\ \hline
  \small
  \textbullet\ Primera herramienta de código abierto en su campo.\newline
  \textbullet\ Transparencia total, permitiendo optimización y adaptabilidad por parte del usuario.\newline & 
  \small
  \textbullet\ Proyecto emergente en un mercado con alternativas consolidadas.\newline
  \textbullet\ Necesidad inicial de cambios para adaptarse a los requisitos de la comunidad.\\ \hline
  \rowcolor{udcpink!25}
  \textbf{Oportunidades} & \textbf{Amenazas} \\ \hline
  \small
  \textbullet\ Carencia de herramientas integradas gratuitas de navegación volumétrica.\newline
  \textbullet\ Gran canal de difusión directa a través de la Asset Store. & 
  \small
  \textbullet\ Sector de los videojuegos en constante cambio y evolución.\newline
  \textbullet\ Crecimiento de motores emergentes competitivos (ej. Godot).\newline
  \textbullet\ Riesgo de aparición de una herramienta oficial integrada en el motor. \\ \hline
  \end{tabular}
  \caption[Matriz DAFO del proyecto]{Matriz \acrshort{dafo} del proyecto}
  \label{tab:dafo}
\end{table}

\subsection{Conclusiones}

A partir de la matriz \acrshort{dafo} expuesta en la Tabla~\ref{tab:dafo}, se concluye que el desarrollo de este sistema es altamente viable.
Si bien el proyecto deberá competir con soluciones comerciales maduras, su naturaleza de código abierto representa una ventaja competitiva crítica.
En definitiva, el proyecto se posiciona como una alternativa democratizadora y accesible, con un claro potencial de adopción dentro del vasto ecosistema de desarrolladores de Unity.
